\documentclass[aspectratio=169]{beamer}
\usetheme{Madrid}
\usecolortheme{default}
\usepackage{amsmath}
\usepackage{physics}

\title{Output Light Analysis}
\subtitle{}
\author{}
\date{\today}

\begin{document}

\frame{\titlepage}

\begin{frame}{Problem}
    From our QuTip simulation we're able to calculate the output field $a_{out}(t)$ via input output theory:
    \begin{equation}
        a_{out}(t) = \mu_{fr}a_{in}(t) + \mu_{fc}\sqrt{\kappa_{oc}}\langle a\rangle(t)
        \label{eq:io-theory}
    \end{equation}
    Where all the fields mentioned in \eqref{eq:io-theory}, can be in one of the following bases: $\{|\pi\rangle,|V\rangle,|+\rangle,|-\rangle \}$

    Now the question becomes, how is one able to analyze this output field and generate a truth table from this. For that we're going to calculate the mean photon number in the output field via:
    $$\bar{n}_{out;\{|\pi\rangle,|V\rangle,|+\rangle,|-\rangle\} } = \int |a_{out;\{ |\pi\rangle,|V\rangle,|+\rangle,|-\rangle  \}}|^2 dt$$

    For the problem we're going to be looking at, we're only going to consider the $|+\rangle$ and the $|-\rangle$ modes of our light, as these are the polarizations which are set in our CNOT basis, for which we want to calculate the fidelity for.
\end{frame}

\begin{frame}{Analysis via detector}
    Now after having extracted the mean photon number in our output field for our two polarizations we need to consider the framework which we're using to analyze this output field. In our case that is the case of 2 independent detectors, where one detects the $|+\rangle$ polarized light and the other detects the $|-\rangle$ polarized light. Both detectors share the same efficiency $\eta$, and also the same dark count rate $p_{dc}$, which is the rate at with a single detector will fire even upon receiving no signal. The efficiency can be viewed as a beam splitter. Meaning, that if we were to have $\bar{n}$ photon in from of our detector that has an efficiency $\eta$, this can be treated as a perfect detector with $\eta=1$ but with a beam splitter in front of our perfect detector that will "take" away $(1-\eta)\bar{n}$ photons from our initial pulse. The rest of the icoming pulse (now: $\eta\bar{n}$) will however get detected perfectly. Since we're only interested in the cases where we have 1 detector event ($k=1$) upon receiving $\lambda=\eta\bar{n}$ photons, we can treat this problem via poisson statistics with:
    $$ P(k,\lambda)=\frac{\lambda^ke^{-\lambda}}{k!} = \lambda e^{-\lambda}$$
    This however doesn't take into consideration our dark counts yet.
\end{frame}

\begin{frame}{Detector statistics}
    If our detector fires, and gives us 1 detector click, it can come from 2 sources:
    \begin{itemize}
        \item a dark count
        \item a photon
    \end{itemize}
    This we can then group up into 4 cases for each of the 2 possible sources of a click
    \begin{itemize}
        \item $P(\text{0 Photons, 0 dark counts}) = e^{-\lambda}\cdot e^{-p_{dc}} = e^{-(\lambda+p_{dc})}$
        \item $P(\text{1 Photon, 0 dark counts}) = \lambda e^{-\lambda}\cdot e^{-p_{dc}} = \lambda e^{-(\lambda+p_{dc})}$
        \item $P(\text{0 Photons, 1 dark count}) = e^{-\lambda}\cdot p_{dc}e^{-p_{dc}} = p_{dc}e^{-(\lambda+p_{pc})}$
        \item $P(\text{1 Photons, 1 dark counts}) =$ Can be ignored as this would produce 2 clicks, which we're post selecting away 
    \end{itemize}

    Meaning that we have 2 out of our 4 possible events which then produce only 1 click, giving us the probability that each of our detectors produces 1 click:
    \begin{equation}
        P_{\pm, click} = (\lambda_\pm + p_{dc})e^{-(\lambda_\pm+p_{dc})}
        \label{eq:prob_click}
    \end{equation}
\end{frame}

\begin{frame}{2 Detector Statistics}
    As mentioned before, we have 2 detectors with the same efficiency $\eta$ and the same detector dark rate $p_{dc}$. These can produce a click independent from one another, but we're only interested in the cases where \textbf{ONLY} 1 produces a click, as the photon can (ideally) only be either $|+\rangle$ or $|-\rangle$ polarized.
    \begin{align*}
        P(\text{only}|+\rangle) &= P_{+,click}\cdot P_{-,0 clicks}\\
                                &= (\lambda_{+} + p_{dc})e^{-(\eta(\bar{n}_{out,+} + \bar{n}_{out,-}) + 2p_{dc})}
    \end{align*}
    \begin{align*}
        P(\text{only}|-\rangle) = (\lambda_{-} + p_{dc})e^{-(\eta(\bar{n}_{out,+} + \bar{n}_{out,-}) + 2p_{dc})}
    \end{align*}

    Adding the two together gives us the probability that at just from the fact that only one detector clicked, something sort of useful came out of our output field, which we're going to call $P_{valid}$:
    $$
    P_{valid} = (\lambda_{+} + \lambda_{-} + 2p_{dc})e^{-(\eta(\bar{n}_{out,+}+\bar{n}_{out,-})+2p_{dc})}
    $$
    However now we need to filter out only the cases where the detector event was actually triggered by a photon and not a dark count.
\end{frame}
\begin{frame}{Photon Event Filtering}
    Here the probability that the 1 detector event in \textbf{ONLY} was caused by a photon is:
    \begin{align*}
        P(\text{1 photon in + \& no click in -}) &= \lambda_{+} e^{-(\lambda_{+} + p_{dc})} \cdot e^{-(\lambda_{-} + p_{dc})}\\
                                                 &= \lambda_{+} e^{-(\eta(\bar{n}_{out,+}+\bar{n}_{out,-})+2p_{dc})} \\
        P(\text{1 photon in - \& no click in +}) &=  \lambda_{-} e^{-(\eta(\bar{n}_{out,+}+\bar{n}_{out,-})+2p_{dc})}
    \end{align*}

    Thus giving us the probability that 1 click was produced in only 1 of the two detectors:
    $$ P_{only} = (\lambda_{+} + \lambda_{-}) e^{-(\eta(\bar{n}_{out,+}+\bar{n}_{out,-})+2p_{dc})}$$

    This then allows us now to calculate how many of our valid clicks, were actually produced by only 1 photon event going only to 1 of the 2 detectors, essentially giving us a sort of signal-to-noise ratio:
    $$
    P_{SNR} = \frac{P_{only}}{P_{valid}} = \frac{\lambda_{+} + \lambda_{-}}{\lambda_{+} + \lambda_{-} + 2p_{dc}}
    $$
\end{frame}
\begin{frame}{Photon Distribution}
    All that's left now is to calculate how much of our light went to the "+"-detector, and how much went to our "-"-detector, which is just like calculating the ratios of all of our events that were triggered only by photons, and the probability of the respective detector itself:
    $$
    R_{+} = \frac{P(\text{1 photon in + and no click in -})}{P_{only}} = \frac{\bar{n}_{out,+}}{\bar{n}_{out,+} + \bar{n}_{out,-}}
    $$
    $$
    R_{-} = \frac{P(\text{1 photon in - and no click in +})}{P_{only}} = \frac{\bar{n}_{out,-}}{\bar{n}_{out,+} + \bar{n}_{out,-}}
    $$

    What then ends up going into our truth table is the following:
    \begin{itemize}
        \item For the $|+\rangle$-part = $R_{+}\cdot P_{SNR} + (1-P_{SNR})\cdot 0.5$
        \item For the $|-\rangle$-part = $R_{-}\cdot P_{SNR} + (1-P_{SNR})\cdot 0.5$
    \end{itemize}
\end{frame}

\end{document}
